\section{Related Work}\label{s:related-work}

\textbf{Sparse Autoencoders.} 
Models can represent more than $n$ features in $n$-dimensional activations through superposition, which complicates interpretability \citep{elhage2022toy}.
Sparse autoencoders (SAEs) \citep{bricken2023monosemanticity, huben2024sparse} and their variants \citep{gao2024scaling, bussmann2024batchtopk, rajamanoharan2024jumping, bussmann2025learning, karvonen2025saebench, leask2025inference} decompose dense activations into sparse feature representations in an attempt to address this.
An SAE consists of a single hidden layer trained to reconstruct model activations subject to a sparsity penalty on the latent activations (typically $L_1$ or top-$k$ regularisation) \citep{huben2024sparse, bricken2023monosemanticity}.

We use 3 types of SAEs in this paper.
JumpReLU SAEs \citep{rajamanoharan2024jumping} use a discontinuous step function at a learned threshold and optimise an $L_0$ penalty.
BatchTopK SAEs \citep{bussmann2024batchtopk} keep the top $n \times k$ activations over a batch of $n$ inputs and sets the rest to zero. This forces an average of $k$ latent activations per input, so the actual $L_0$ sparsity is approximately $k$. 
Matryoshka SAEs \citep{bussmann2025learning} train nested sub-dictionaries of increasing size under a shared encoder and decoder, producing representations at multiple levels of granularity from a single training run.
We focus on BatchTopK SAEs because they are much simpler than the others.
% TODO justify this? or do in method? i saw a paper talking about using btk for simplicity or smthn but can't remember it!

SAEs have been successfully used in the past \citep{marks2025auditing, movva2025sparse} but there are uncertainties over whether they learn robust dictionaries \citep{chanin2024absorption,leask2025sparse}. In our paper, we show that even a correct dictionary may be insufficient for complete model understanding using SAEs.


\textbf{Relative Magnitude-based Composition.}
\cite{farrell2025order} train a two-layer attention-only transformer to solve an ordered bigram copying task, in which the model must encode bigram $(d_1,d_2)$ into a single separator token and then reconstruct the bigram in the same order as output.
They find that the model learns to represent the ordered bigram at the separator token using the relative magnitude of scalar attention weights, rather than representing each possible ordering as a different direction in activation space.
Specifically, they calculate the residual activation at the separator token $\boldsymbol r^{L_1}_s$ as:
\begin{equation*}
    \boldsymbol r^{L_1}_s = \alpha_{s\to d_1} (\boldsymbol E_{d_1} + \boldsymbol P_{d_1}) + \alpha_{s \to d_2}(\boldsymbol E_{d_2} + \boldsymbol P_{d_2}) + \boldsymbol E_s + \boldsymbol P_s
\end{equation*}
where $\alpha_{s\to d_i}$ is the attention weight in the first layer for query $s$ (separator token) and key $d_i$ (element of input bigram), and $\boldsymbol E_{x}$ and $\boldsymbol P_{x}$ are the token and positional embeddings for token $x$ respectively.
They predict that an SAE trained on $\boldsymbol r^{L_1}_s$ will learn latents corresponding to $(\boldsymbol E_{a} + \boldsymbol P_{d_1})$ and $(\boldsymbol E_{b} + \boldsymbol P_{d_2})$ with activations equal to the attention probabilities $\alpha$, and that this results in graded latent activations.

Graded latent activations undermine the assumption that SAE latents are independent and binary \cite{} because ...
% which would mean latents exhibit different behavior at different activation magnitudes, rather than simply being on/off

\textbf{Steering.}
TODO


% \textbf{Mechanistic Interpretability.}
% Mechanistic interpretability is the study of reverse-engineering neural networks 
% into human-interpretable algorithms \citep{olah2020zoom, 
% ferrando2024primer, nanda2023progress}. This is usually achieved by identifying ``circuits'' in the model: paths of computation through a model that are responsible for specific behaviours, represented as model subgraphs.
% Some example behaviours that have been successfully reverse-engineered in this way are indirect object identification \citep{wang2022interpretability}, in-context learning via induction heads \citep{olsson2022context}, and modular arithmetic via Fourier representations \citep{nanda2023progress}. 
% Many of these studies use toy models as a controlled environment for identifying mechanisms that might generalise to larger models, which supports our use of toy models in this paper. For example, \citet{baeumel2025modular} found the modular arithmetic results on toy models subsequently generalised to pretrained LLMs.

% Moreover, recent work has found ``steering'' vectors in the residual stream which can be steered via linear scaling and patched back into the residual vector to predictably change model behaviour. For example,
% \citet{arditi2024refusal} use the difference-in-means vector between residual streams on harmful and harmless instructions to identify a single ``refusal direction'' in the residual stream of safety-tuned LMs, and use it to to suppress or amplify refusal behaviour via linear scaling.
% % adding negative multiples of a direction correlated with a given output sentiment ($\boldsymbol u$) to the residual stream ($\boldsymbol r{L\_i}' \leftarrow \boldsymbol r{L\_i} - \alpha \boldsymbol u$) can cause the model to generate output with the opposite sentiment \cite{tigges2024language}
% This provides evidence that a given direction in the residual stream is causally relevant for a given behaviour, which provides insights into what feature that direction represents \cite{ferrando2024primer}.
% We use the same principle in our latent steering experiments, in which we steer SAE latents rather than raw residual stream directions and patch the SAE-reconstructed activations into the residual stream.

% \textbf{Relational composition and binding.}

% \begin{table*}[t]
% \centering
% \caption{Comparison of relational composition mechanisms and their implications for interpretability.}
% \label{tab:composition_mechanisms}
% \begin{tabular}{p{3.5cm}p{4cm}p{4cm}p{4cm}}
% \toprule
% \textbf{Composition} & \textbf{Description} & \textbf{Example} & \textbf{Implications} \\
% \midrule
% Direction-based &
% Features are vectors; sequence position defined by different directions.
% & Matrix binding \cite{wattenberg2024relational} - features projected to distinct linear subspaces: $r = Ax + By$. &
% Feature multiplicity: SAEs may learn echo features ($Ax$, $Ay$, $Bx$, $By$) representing the same underlying concept. \\
% \addlinespace
% Fixed magnitude-based &
% Multiple features have same direction; sequence position defined by different magnitudes.
% &
% Onion-like representations \cite{csordas-etal-2024-recurrent} - ``unpeel'' via autoregression to retrieve sequence positions.
% &
% Counter-example to the widely held assumption that representations are linear.
% \\
% \addlinespace
% Relative magnitude-based &
% Sequence position defined by magnitude relative to other positions in sequence.
% & This work - sequence order defined by a weighted sum of input token and positional embeddings. &
% Counter-example to the widely held assumptions that representations are linear and SAE latents are binary.
% \\
% \bottomrule
% \end{tabular}
% \end{table*}

% Classic work in distributed representations show how single vectors can encode structured data. Tensor Product Representations (TRPs) and Vector Symbolic Architectures (VSAs) bind and additively superpose vectors, enabling ordered structures in fixed width vectors \citep{smolensky1990tpr, plate1995hrr, kanerva2009hdc}.

% \citet{wattenberg2024relational} highlight additive matrix binding, where vectors are written to slots using role-specific linear maps so that bigrams $(x, y),(y,x)$ remain separable, as a candidate mechanism for relational composition in transformers. That is, given an ordered bigram $(d_1, d_2)$ and two $n \times n$ matrices $A$ and $B$, define the representation of $(x, y)$ by
% \(
% r = Ax + By
% \).
% In this way, $(x, y)$ and $(y, x)$ have different representations.
% \citet{wattenberg2024relational} claim that this causes \textit{feature multiplicity}: a kind of false positive where multiple ``echo features'' are created that represent the same concept in different contexts. Given two features $x$ and $y$, matrix binding results in $Ax$, $Ay$, $Bx$, and $By$ features that represent the same concept as $x$ and $y$. 

% Empirically, \citet{feng2023language, prakash2023fine} identify binding ID vectors that bind entities to their attributes in-context in language models. For example, a ``green car'' is represented by a vector representing ``green'', a vector for ``car'', and a binding ID vector to join the two. A binding ID vector $k$ is attached to an entity and its attribute by addition in their respective activation spaces. To answer a query for a given entity, the attribute with a matching binding ID is retrieved from the context activation space. \citet{feng2023language} also show that these mechanisms are position-independent with respect to the attribute and entity, and that the activations can be factorised.

% \citet{brinkmann2024mechanistic} reverse-engineer an attention-only tree-search transformer that stores precomputed subpaths in special token positions and merges them later. This occurs when the tree depth, $N$, is greater than $L-1$, where $L$ is the number of transformer layers. This constraint is also necessary in our work, where $N$ is the size of the input $N$-gram. However, the authors do not investigate the structure of these activations.

% \citet{csordas-etal-2024-recurrent} demonstrate ``onion representations'' in small gated recurrent neural networks (RNNs), where these slots are defined by different orders of magnitude rather than distinct linear subspaces.
% The model represents ordered sequences by closing the gates gradually and synchronously over the input phase, which exponentially decays the scaling factor of subsequent sequence positions.
% This produces layered onion representations where earlier sequence positions dominate the magnitude space.
% This contrasts larger models which indicate sequence position by sharply closing their gates, which creates position-dependent subspaces for each input.
% The onion representation allows autoregressive decoding to sequentially ``peel'' off tokens by subtracting the current dominating token embedding.
% Multiple tokens can occupy the same directional subspace at different magnitude scales; any linear direction will cross-cut multiple layers of the onion. 
% While the authors find that sequence order is represented by fixed magnitude scales, we find an alternative mechanism where it is instead represented by relative magnitude between positions, which is input-dependent and not fixed.

% Our toy model provides evidence for magnitude-based composition in attention-only transformers, but with \textit{relative} magnitudes rather than fixed ones as observed in \citet{csordas-etal-2024-recurrent}. 
% % Rather than each position being assigned a predetermined magnitude scale $\gamma^t$, the compositional information is encoded in the relative scaling between feature components that can adapt to the specific features being combined. Additionally, unlike GRUs which require autoregressive unpeeling to decode onion representations, transformers can potentially access multiple positions in parallel through attention mechanisms.
% See Table \ref{tab:composition_mechanisms} for an overview of these composition methods.
